PhysAgent combines a TCAD physics ontology and literature-driven domain knowledge to guide LLM-based SDevice file generation. The constructed ontology is publicly available as an interactive web platform~\cite{ontology_tcad}.

\subsection{TCAD Physics Ontology}
An SDevice simulation takes a command file (\texttt{.cmd}) and a parameter file (\texttt{.par}) as input. The physics block in the \texttt{.cmd} file specifies the physical models activated for each material, while the \texttt{.par} file defines the parameters required by these models along with fundamental properties such as bandgap, permittivity, and the DOS effective masses of electrons and holes. Thus, the model selection in the \texttt{.cmd} file determines the parameter set that must be provided in the \texttt{.par} file.
To consistently capture this \texttt{.cmd}--\texttt{.par} dependency, we structure the relevant parameters, TCAD syntax, and physical knowledge into an ontology grounded in {\color{red}the Poisson and transport equations, together with the governing equations of each physical model}. {\color{red}Each parameter and each physical-model equation is represented as a node. Parameter nodes attach to the equation node of the model they belong to (membership), and equation nodes link to one another to capture inter-model couplings (coupling). Given a model selection, traversing these edges yields the required parameter set.} A syntax layer then applies the SDevice syntax to the selected equations and nodes to generate the \texttt{.cmd} and \texttt{.par} files (Fig.~\ref{fig:ontology}).

\subsection{Ontology-Guided Generation}
Given a target device and material, PhysAgent accesses the ontology through an MCP (Model Context Protocol) interface {\color{red}(referred to as \texttt{ontology\_mcp})} that exposes four tools. \texttt{ontology\_overview} returns the node/edge map of the graph, serving as the entry point for grasping the overall structure. \texttt{ontology\_get\_chapter} identifies the governing equations associated with a selected node and derives the parameter set required by those equations. \texttt{ontology\_get\_syntax} provides the SDevice syntax template for each model, specifying keyword names, section layout, and required fields. \texttt{ontology\_get\_reference} returns material-specific parameter values together with per-parameter source citations drawn from a curated collection of papers, datasheets, and TCAD references{\color{red}~\cite{bludau1974, ioffe_nsm}}. In our experiments, we employ Claude Sonnet, whose context window is limited to 200k tokens, whereas the full ontology measures $\sim$458k tokens and therefore cannot be injected directly into the prompt. The MCP-based selective retrieval scheme circumvents this limitation by allowing the LLM to query only the information required at each step, keeping the per-query context well within the 200k-token budget. Moreover, each parameter in the generated \texttt{.par} inherits the citation returned by \texttt{ontology\_get\_reference}, thereby enabling parameter-level provenance (Fig.~\ref{fig:exp2}).

